% ============================================================================
%  overlay-core.tex —— 投標表單套印的共用巨集
%
%  概念：用 pdfpages 載入機關提供的原始表單 PDF，只在欄位位置「疊印」廠商資料，
%        原始表單完全不更動。座標單位 bp（PDF 點），原點在頁面左下角，y 是文字基線。
%
%  使用方式：在你的 forms.tex 裡
%      \input{vendor.tex}            % 先載入資料
%      % ============================================================================
%  overlay-core.tex —— 投標表單套印的共用巨集
%
%  概念：用 pdfpages 載入機關提供的原始表單 PDF，只在欄位位置「疊印」廠商資料，
%        原始表單完全不更動。座標單位 bp（PDF 點），原點在頁面左下角，y 是文字基線。
%
%  使用方式：在你的 forms.tex 裡
%      \input{vendor.tex}            % 先載入資料
%      % ============================================================================
%  overlay-core.tex —— 投標表單套印的共用巨集
%
%  概念：用 pdfpages 載入機關提供的原始表單 PDF，只在欄位位置「疊印」廠商資料，
%        原始表單完全不更動。座標單位 bp（PDF 點），原點在頁面左下角，y 是文字基線。
%
%  使用方式：在你的 forms.tex 裡
%      \input{vendor.tex}            % 先載入資料
%      % ============================================================================
%  overlay-core.tex —— 投標表單套印的共用巨集
%
%  概念：用 pdfpages 載入機關提供的原始表單 PDF，只在欄位位置「疊印」廠商資料，
%        原始表單完全不更動。座標單位 bp（PDF 點），原點在頁面左下角，y 是文字基線。
%
%  使用方式：在你的 forms.tex 裡
%      \input{vendor.tex}            % 先載入資料
%      \input{../core/overlay-core}  % 再載入本檔
%  然後用 \includepdf[..., picturecommand*={\你的Overlay}]{src/表單.pdf}
%
%  值的三種狀態（填在 vendor.tex）：
%      {}        空白      → 草稿印紅色〔提示〕；正式版不印
%      {否?}     結尾有 ?  → 建議值、尚未確認：草稿印紅字；預覽版印黑字；正式版不印
%      {否}      其他      → 已確認：一律印黑字
%  模式：\drafttrue（草稿）／\draftfalse（正式）；preview.tex 會設 \PreviewMode（預覽版）。
% ============================================================================
\usepackage{fontspec}
\usepackage{xeCJK}
\usepackage{xcolor}
\usepackage{graphicx}
\usepackage{etoolbox}
\usepackage{xstring}
\expandarg
\usepackage{pdfpages}

% ---------------------------------------------------------------- 字型
% 依序找得到哪個就用哪個；要指定可在 vendor.tex 裡先 \def\TenderCJKFont{字型名稱}
\providecommand{\TenderCJKFont}{}
\providecommand{\TenderLatinFont}{}
\newcommand{\tryCJK}[1]{\ifdefempty{\TenderCJKFont}{\IfFontExistsTF{#1}{\def\TenderCJKFont{#1}}{}}{}}
\tryCJK{DFKai-SB}               % Windows 標楷體
\tryCJK{BiauKai}                % macOS 標楷體
\tryCJK{TW-Kai}                 % 全字庫正楷體
\tryCJK{Noto Serif CJK TC}
\tryCJK{Noto Sans CJK TC}
\ifdefempty{\TenderCJKFont}{\PackageError{overlay-core}{找不到中文字型，請在 vendor.tex 設定 \string\TenderCJKFont}{}}{}
\setCJKmainfont{\TenderCJKFont}
\ifdefempty{\TenderLatinFont}{%
  \IfFontExistsTF{Times New Roman}{\setmainfont{Times New Roman}}{\setmainfont{TeX Gyre Termes}}}%
  {\setmainfont{\TenderLatinFont}}

% \ifdraft、\ifpreview 由 forms.tex 在載入 vendor.tex 之前用 \newif 宣告
\ifdefined\PreviewMode\draftfalse\previewtrue\fi
\setlength{\unitlength}{1bp}

% ---------------------------------------------------------------- 放字
% \putl / \putc / \putr [字級]{x,y}{內容}：靠左／置中／靠右，y 為基線
\newcommand{\putl}[3][12]{\put(#2){\fontsize{#1}{#1}\selectfont #3}}
\newcommand{\putc}[3][12]{\put(#2){\makebox[0pt]{\fontsize{#1}{#1}\selectfont #3}}}
\newcommand{\putr}[3][12]{\put(#2){\makebox[0pt][r]{\fontsize{#1}{#1}\selectfont #3}}}

% ---------------------------------------------------------------- 確認狀態
% \ifunconf{值}{是}{否}：值是否以 ? 或 ？ 結尾
\newcommand{\ifunconf}[3]{\IfEndWith{#1}{?}{#2}{\IfEndWith{#1}{？}{#2}{#3}}}
% \unc{文字}：尚未確認的值——草稿紅字、預覽版黑字、正式版不印
\newcommand{\unc}[1]{\ifdraft\textcolor{red}{#1}\fi\ifpreview#1\fi}
% \val{\巨集}{提示}：必填欄位（空白時草稿印提示）
\newcommand{\val}[2]{\ifdefempty{#1}{\ifdraft\textcolor{red}{〔#2〕}\fi}%
  {\ifunconf{#1}{\unc{\StrGobbleRight{#1}{1}}}{#1}}}
% \opt{\巨集}：選填欄位（空白時不印）
\newcommand{\opt}[1]{\ifdefempty{#1}{}{\ifunconf{#1}{\unc{\StrGobbleRight{#1}{1}}}{#1}}}
% \draftnote{y}：頁首狀態說明（草稿紅字／預覽版灰字）
\newcommand{\draftnote}[1]{%
  \ifdraft\putl[8]{28,#1}{\textcolor{red}{草稿｜紅字＝尚未填寫或尚未確認。在 vendor.tex 補資料；值後面的「?」刪掉即視為確認。}}\fi
  \ifpreview\putl[8]{28,#1}{\textcolor{gray}{預覽版｜帶「?」的建議值也以黑字顯示，僅供檢查版面，勿用於投標。}}\fi}

% ---------------------------------------------------------------- 版面輔助
% \fitw{最大寬度}{內容}：超過寬度就等比縮小（長公司名、地址）
\newsavebox{\fitbox}
\newcommand{\fitw}[2]{\sbox{\fitbox}{#2}%
  \ifdim\wd\fitbox>#1\relax\resizebox{#1}{!}{\usebox{\fitbox}}\else\usebox{\fitbox}\fi}
% \putwrap[字級]{x,y}{寬度}{最大高度}{內容}：以 x 置中；放不下自動換行，再放不下等比縮小
\newsavebox{\wrapA}\newsavebox{\wrapB}
\newcommand{\putwrap}[5][14]{%
  \sbox{\wrapA}{\fontsize{#1}{#1}\selectfont #5}%
  \ifdim\wd\wrapA>#3\relax
    \sbox{\wrapA}{\parbox{#3}{\centering\fontsize{#1}{\dimexpr#1pt*23/20\relax}\selectfont #5}}%
    \sbox{\wrapB}{\ifdim\dimexpr\ht\wrapA+\dp\wrapA\relax>#4\relax\resizebox*{!}{#4}{\usebox{\wrapA}}\else\usebox{\wrapA}\fi}%
    \put(#2){\makebox[0pt]{\raisebox{\dimexpr#1pt*35/100-\ht\wrapB/2+\dp\wrapB/2\relax}{\usebox{\wrapB}}}}%
  \else
    \put(#2){\makebox[0pt]{\usebox{\wrapA}}}%
  \fi}

% ---------------------------------------------------------------- 勾選
\newcommand{\vmark}[2][13]{\put(#2){\makebox[0pt]{\fontsize{#1}{#1}\selectfont\bfseries V}}}
\newcommand{\rvmark}[2][13]{\put(#2){\makebox[0pt]{\fontsize{#1}{#1}\selectfont\bfseries\textcolor{red}{V}}}}
\newcommand{\uvmark}[2][13]{\ifdraft\rvmark[#1]{#2}\fi\ifpreview\vmark[#1]{#2}\fi}
% \chk{\開關}{x,y}：1 = 勾；1? = 尚未確認的勾；其他不勾
\newcommand{\chk}[2]{\ifdefstring{#1}{1}{\vmark[11]{#2}}{%
  \ifdefstring{#1}{1?}{\uvmark[11]{#2}}{\ifdefstring{#1}{1？}{\uvmark[11]{#2}}{}}}}

% \decl[V 字級]{\巨集}{y}：是／否兩欄的勾選（聲明書這類表格）。
%   欄位中心 x 由 \DeclYesX、\DeclNoX 設定（每份表單量一次）。
\providecommand{\DeclYesX}{474}
\providecommand{\DeclNoX}{527.5}
\newcommand{\decl}[3][13]{%
  \ifdefstring{#2}{是}{\vmark[#1]{\DeclYesX,#3}}{%
  \ifdefstring{#2}{否}{\vmark[#1]{\DeclNoX,#3}}{%
  \ifdefstring{#2}{是?}{\uvmark[#1]{\DeclYesX,#3}}{%
  \ifdefstring{#2}{否?}{\uvmark[#1]{\DeclNoX,#3}}{%
  \ifdefstring{#2}{是？}{\uvmark[#1]{\DeclYesX,#3}}{%
  \ifdefstring{#2}{否？}{\uvmark[#1]{\DeclNoX,#3}}{%
  \ifdraft\put(\DeclYesX,#3){\makebox[\dimexpr\DeclNoX bp-\DeclYesX bp\relax]{%
    \fontsize{9}{9}\selectfont\textcolor{red}{〔待定〕}}}\fi}}}}}}}

% \digit{\巨集}{第幾個字}{x,y}{字級}：把字串第 n 個字置中印在格子裡（中文大寫金額、郵遞區號）
\newcommand{\digit}[4]{\putc[#4]{#3}{%
  \ifunconf{#1}{\StrGobbleRight{#1}{1}[\digitstr]\unc{\StrChar{\digitstr}{#2}}}{\StrChar{#1}{#2}}}}

% ---------------------------------------------------------------- 開關
% \repname[字級]{x,y}{姓名}：\PrintRepOnSign = 1 才在簽章處印負責人姓名
\providecommand{\PrintRepOnSign}{0}
\newcommand{\repname}[3][12]{\ifdefstring{\PrintRepOnSign}{1}{\putl[#1]{#2}{#3}}{}}
% \ifsigndate{內容}：\PrintSignDate = 1 才印簽署日期（0 = 留白手寫）
\providecommand{\PrintSignDate}{1}
\newcommand{\ifsigndate}[1]{\ifdefstring{\PrintSignDate}{1}{#1}{}}
  % 再載入本檔
%  然後用 \includepdf[..., picturecommand*={\你的Overlay}]{src/表單.pdf}
%
%  值的三種狀態（填在 vendor.tex）：
%      {}        空白      → 草稿印紅色〔提示〕；正式版不印
%      {否?}     結尾有 ?  → 建議值、尚未確認：草稿印紅字；預覽版印黑字；正式版不印
%      {否}      其他      → 已確認：一律印黑字
%  模式：\drafttrue（草稿）／\draftfalse（正式）；preview.tex 會設 \PreviewMode（預覽版）。
% ============================================================================
\usepackage{fontspec}
\usepackage{xeCJK}
\usepackage{xcolor}
\usepackage{graphicx}
\usepackage{etoolbox}
\usepackage{xstring}
\expandarg
\usepackage{pdfpages}

% ---------------------------------------------------------------- 字型
% 依序找得到哪個就用哪個；要指定可在 vendor.tex 裡先 \def\TenderCJKFont{字型名稱}
\providecommand{\TenderCJKFont}{}
\providecommand{\TenderLatinFont}{}
\newcommand{\tryCJK}[1]{\ifdefempty{\TenderCJKFont}{\IfFontExistsTF{#1}{\def\TenderCJKFont{#1}}{}}{}}
\tryCJK{DFKai-SB}               % Windows 標楷體
\tryCJK{BiauKai}                % macOS 標楷體
\tryCJK{TW-Kai}                 % 全字庫正楷體
\tryCJK{Noto Serif CJK TC}
\tryCJK{Noto Sans CJK TC}
\ifdefempty{\TenderCJKFont}{\PackageError{overlay-core}{找不到中文字型，請在 vendor.tex 設定 \string\TenderCJKFont}{}}{}
\setCJKmainfont{\TenderCJKFont}
\ifdefempty{\TenderLatinFont}{%
  \IfFontExistsTF{Times New Roman}{\setmainfont{Times New Roman}}{\setmainfont{TeX Gyre Termes}}}%
  {\setmainfont{\TenderLatinFont}}

% \ifdraft、\ifpreview 由 forms.tex 在載入 vendor.tex 之前用 \newif 宣告
\ifdefined\PreviewMode\draftfalse\previewtrue\fi
\setlength{\unitlength}{1bp}

% ---------------------------------------------------------------- 放字
% \putl / \putc / \putr [字級]{x,y}{內容}：靠左／置中／靠右，y 為基線
\newcommand{\putl}[3][12]{\put(#2){\fontsize{#1}{#1}\selectfont #3}}
\newcommand{\putc}[3][12]{\put(#2){\makebox[0pt]{\fontsize{#1}{#1}\selectfont #3}}}
\newcommand{\putr}[3][12]{\put(#2){\makebox[0pt][r]{\fontsize{#1}{#1}\selectfont #3}}}

% ---------------------------------------------------------------- 確認狀態
% \ifunconf{值}{是}{否}：值是否以 ? 或 ？ 結尾
\newcommand{\ifunconf}[3]{\IfEndWith{#1}{?}{#2}{\IfEndWith{#1}{？}{#2}{#3}}}
% \unc{文字}：尚未確認的值——草稿紅字、預覽版黑字、正式版不印
\newcommand{\unc}[1]{\ifdraft\textcolor{red}{#1}\fi\ifpreview#1\fi}
% \val{\巨集}{提示}：必填欄位（空白時草稿印提示）
\newcommand{\val}[2]{\ifdefempty{#1}{\ifdraft\textcolor{red}{〔#2〕}\fi}%
  {\ifunconf{#1}{\unc{\StrGobbleRight{#1}{1}}}{#1}}}
% \opt{\巨集}：選填欄位（空白時不印）
\newcommand{\opt}[1]{\ifdefempty{#1}{}{\ifunconf{#1}{\unc{\StrGobbleRight{#1}{1}}}{#1}}}
% \draftnote{y}：頁首狀態說明（草稿紅字／預覽版灰字）
\newcommand{\draftnote}[1]{%
  \ifdraft\putl[8]{28,#1}{\textcolor{red}{草稿｜紅字＝尚未填寫或尚未確認。在 vendor.tex 補資料；值後面的「?」刪掉即視為確認。}}\fi
  \ifpreview\putl[8]{28,#1}{\textcolor{gray}{預覽版｜帶「?」的建議值也以黑字顯示，僅供檢查版面，勿用於投標。}}\fi}

% ---------------------------------------------------------------- 版面輔助
% \fitw{最大寬度}{內容}：超過寬度就等比縮小（長公司名、地址）
\newsavebox{\fitbox}
\newcommand{\fitw}[2]{\sbox{\fitbox}{#2}%
  \ifdim\wd\fitbox>#1\relax\resizebox{#1}{!}{\usebox{\fitbox}}\else\usebox{\fitbox}\fi}
% \putwrap[字級]{x,y}{寬度}{最大高度}{內容}：以 x 置中；放不下自動換行，再放不下等比縮小
\newsavebox{\wrapA}\newsavebox{\wrapB}
\newcommand{\putwrap}[5][14]{%
  \sbox{\wrapA}{\fontsize{#1}{#1}\selectfont #5}%
  \ifdim\wd\wrapA>#3\relax
    \sbox{\wrapA}{\parbox{#3}{\centering\fontsize{#1}{\dimexpr#1pt*23/20\relax}\selectfont #5}}%
    \sbox{\wrapB}{\ifdim\dimexpr\ht\wrapA+\dp\wrapA\relax>#4\relax\resizebox*{!}{#4}{\usebox{\wrapA}}\else\usebox{\wrapA}\fi}%
    \put(#2){\makebox[0pt]{\raisebox{\dimexpr#1pt*35/100-\ht\wrapB/2+\dp\wrapB/2\relax}{\usebox{\wrapB}}}}%
  \else
    \put(#2){\makebox[0pt]{\usebox{\wrapA}}}%
  \fi}

% ---------------------------------------------------------------- 勾選
\newcommand{\vmark}[2][13]{\put(#2){\makebox[0pt]{\fontsize{#1}{#1}\selectfont\bfseries V}}}
\newcommand{\rvmark}[2][13]{\put(#2){\makebox[0pt]{\fontsize{#1}{#1}\selectfont\bfseries\textcolor{red}{V}}}}
\newcommand{\uvmark}[2][13]{\ifdraft\rvmark[#1]{#2}\fi\ifpreview\vmark[#1]{#2}\fi}
% \chk{\開關}{x,y}：1 = 勾；1? = 尚未確認的勾；其他不勾
\newcommand{\chk}[2]{\ifdefstring{#1}{1}{\vmark[11]{#2}}{%
  \ifdefstring{#1}{1?}{\uvmark[11]{#2}}{\ifdefstring{#1}{1？}{\uvmark[11]{#2}}{}}}}

% \decl[V 字級]{\巨集}{y}：是／否兩欄的勾選（聲明書這類表格）。
%   欄位中心 x 由 \DeclYesX、\DeclNoX 設定（每份表單量一次）。
\providecommand{\DeclYesX}{474}
\providecommand{\DeclNoX}{527.5}
\newcommand{\decl}[3][13]{%
  \ifdefstring{#2}{是}{\vmark[#1]{\DeclYesX,#3}}{%
  \ifdefstring{#2}{否}{\vmark[#1]{\DeclNoX,#3}}{%
  \ifdefstring{#2}{是?}{\uvmark[#1]{\DeclYesX,#3}}{%
  \ifdefstring{#2}{否?}{\uvmark[#1]{\DeclNoX,#3}}{%
  \ifdefstring{#2}{是？}{\uvmark[#1]{\DeclYesX,#3}}{%
  \ifdefstring{#2}{否？}{\uvmark[#1]{\DeclNoX,#3}}{%
  \ifdraft\put(\DeclYesX,#3){\makebox[\dimexpr\DeclNoX bp-\DeclYesX bp\relax]{%
    \fontsize{9}{9}\selectfont\textcolor{red}{〔待定〕}}}\fi}}}}}}}

% \digit{\巨集}{第幾個字}{x,y}{字級}：把字串第 n 個字置中印在格子裡（中文大寫金額、郵遞區號）
\newcommand{\digit}[4]{\putc[#4]{#3}{%
  \ifunconf{#1}{\StrGobbleRight{#1}{1}[\digitstr]\unc{\StrChar{\digitstr}{#2}}}{\StrChar{#1}{#2}}}}

% ---------------------------------------------------------------- 開關
% \repname[字級]{x,y}{姓名}：\PrintRepOnSign = 1 才在簽章處印負責人姓名
\providecommand{\PrintRepOnSign}{0}
\newcommand{\repname}[3][12]{\ifdefstring{\PrintRepOnSign}{1}{\putl[#1]{#2}{#3}}{}}
% \ifsigndate{內容}：\PrintSignDate = 1 才印簽署日期（0 = 留白手寫）
\providecommand{\PrintSignDate}{1}
\newcommand{\ifsigndate}[1]{\ifdefstring{\PrintSignDate}{1}{#1}{}}
  % 再載入本檔
%  然後用 \includepdf[..., picturecommand*={\你的Overlay}]{src/表單.pdf}
%
%  值的三種狀態（填在 vendor.tex）：
%      {}        空白      → 草稿印紅色〔提示〕；正式版不印
%      {否?}     結尾有 ?  → 建議值、尚未確認：草稿印紅字；預覽版印黑字；正式版不印
%      {否}      其他      → 已確認：一律印黑字
%  模式：\drafttrue（草稿）／\draftfalse（正式）；preview.tex 會設 \PreviewMode（預覽版）。
% ============================================================================
\usepackage{fontspec}
\usepackage{xeCJK}
\usepackage{xcolor}
\usepackage{graphicx}
\usepackage{etoolbox}
\usepackage{xstring}
\expandarg
\usepackage{pdfpages}

% ---------------------------------------------------------------- 字型
% 依序找得到哪個就用哪個；要指定可在 vendor.tex 裡先 \def\TenderCJKFont{字型名稱}
\providecommand{\TenderCJKFont}{}
\providecommand{\TenderLatinFont}{}
\newcommand{\tryCJK}[1]{\ifdefempty{\TenderCJKFont}{\IfFontExistsTF{#1}{\def\TenderCJKFont{#1}}{}}{}}
\tryCJK{DFKai-SB}               % Windows 標楷體
\tryCJK{BiauKai}                % macOS 標楷體
\tryCJK{TW-Kai}                 % 全字庫正楷體
\tryCJK{Noto Serif CJK TC}
\tryCJK{Noto Sans CJK TC}
\ifdefempty{\TenderCJKFont}{\PackageError{overlay-core}{找不到中文字型，請在 vendor.tex 設定 \string\TenderCJKFont}{}}{}
\setCJKmainfont{\TenderCJKFont}
\ifdefempty{\TenderLatinFont}{%
  \IfFontExistsTF{Times New Roman}{\setmainfont{Times New Roman}}{\setmainfont{TeX Gyre Termes}}}%
  {\setmainfont{\TenderLatinFont}}

% \ifdraft、\ifpreview 由 forms.tex 在載入 vendor.tex 之前用 \newif 宣告
\ifdefined\PreviewMode\draftfalse\previewtrue\fi
\setlength{\unitlength}{1bp}

% ---------------------------------------------------------------- 放字
% \putl / \putc / \putr [字級]{x,y}{內容}：靠左／置中／靠右，y 為基線
\newcommand{\putl}[3][12]{\put(#2){\fontsize{#1}{#1}\selectfont #3}}
\newcommand{\putc}[3][12]{\put(#2){\makebox[0pt]{\fontsize{#1}{#1}\selectfont #3}}}
\newcommand{\putr}[3][12]{\put(#2){\makebox[0pt][r]{\fontsize{#1}{#1}\selectfont #3}}}

% ---------------------------------------------------------------- 確認狀態
% \ifunconf{值}{是}{否}：值是否以 ? 或 ？ 結尾
\newcommand{\ifunconf}[3]{\IfEndWith{#1}{?}{#2}{\IfEndWith{#1}{？}{#2}{#3}}}
% \unc{文字}：尚未確認的值——草稿紅字、預覽版黑字、正式版不印
\newcommand{\unc}[1]{\ifdraft\textcolor{red}{#1}\fi\ifpreview#1\fi}
% \val{\巨集}{提示}：必填欄位（空白時草稿印提示）
\newcommand{\val}[2]{\ifdefempty{#1}{\ifdraft\textcolor{red}{〔#2〕}\fi}%
  {\ifunconf{#1}{\unc{\StrGobbleRight{#1}{1}}}{#1}}}
% \opt{\巨集}：選填欄位（空白時不印）
\newcommand{\opt}[1]{\ifdefempty{#1}{}{\ifunconf{#1}{\unc{\StrGobbleRight{#1}{1}}}{#1}}}
% \draftnote{y}：頁首狀態說明（草稿紅字／預覽版灰字）
\newcommand{\draftnote}[1]{%
  \ifdraft\putl[8]{28,#1}{\textcolor{red}{草稿｜紅字＝尚未填寫或尚未確認。在 vendor.tex 補資料；值後面的「?」刪掉即視為確認。}}\fi
  \ifpreview\putl[8]{28,#1}{\textcolor{gray}{預覽版｜帶「?」的建議值也以黑字顯示，僅供檢查版面，勿用於投標。}}\fi}

% ---------------------------------------------------------------- 版面輔助
% \fitw{最大寬度}{內容}：超過寬度就等比縮小（長公司名、地址）
\newsavebox{\fitbox}
\newcommand{\fitw}[2]{\sbox{\fitbox}{#2}%
  \ifdim\wd\fitbox>#1\relax\resizebox{#1}{!}{\usebox{\fitbox}}\else\usebox{\fitbox}\fi}
% \putwrap[字級]{x,y}{寬度}{最大高度}{內容}：以 x 置中；放不下自動換行，再放不下等比縮小
\newsavebox{\wrapA}\newsavebox{\wrapB}
\newcommand{\putwrap}[5][14]{%
  \sbox{\wrapA}{\fontsize{#1}{#1}\selectfont #5}%
  \ifdim\wd\wrapA>#3\relax
    \sbox{\wrapA}{\parbox{#3}{\centering\fontsize{#1}{\dimexpr#1pt*23/20\relax}\selectfont #5}}%
    \sbox{\wrapB}{\ifdim\dimexpr\ht\wrapA+\dp\wrapA\relax>#4\relax\resizebox*{!}{#4}{\usebox{\wrapA}}\else\usebox{\wrapA}\fi}%
    \put(#2){\makebox[0pt]{\raisebox{\dimexpr#1pt*35/100-\ht\wrapB/2+\dp\wrapB/2\relax}{\usebox{\wrapB}}}}%
  \else
    \put(#2){\makebox[0pt]{\usebox{\wrapA}}}%
  \fi}

% ---------------------------------------------------------------- 勾選
\newcommand{\vmark}[2][13]{\put(#2){\makebox[0pt]{\fontsize{#1}{#1}\selectfont\bfseries V}}}
\newcommand{\rvmark}[2][13]{\put(#2){\makebox[0pt]{\fontsize{#1}{#1}\selectfont\bfseries\textcolor{red}{V}}}}
\newcommand{\uvmark}[2][13]{\ifdraft\rvmark[#1]{#2}\fi\ifpreview\vmark[#1]{#2}\fi}
% \chk{\開關}{x,y}：1 = 勾；1? = 尚未確認的勾；其他不勾
\newcommand{\chk}[2]{\ifdefstring{#1}{1}{\vmark[11]{#2}}{%
  \ifdefstring{#1}{1?}{\uvmark[11]{#2}}{\ifdefstring{#1}{1？}{\uvmark[11]{#2}}{}}}}

% \decl[V 字級]{\巨集}{y}：是／否兩欄的勾選（聲明書這類表格）。
%   欄位中心 x 由 \DeclYesX、\DeclNoX 設定（每份表單量一次）。
\providecommand{\DeclYesX}{474}
\providecommand{\DeclNoX}{527.5}
\newcommand{\decl}[3][13]{%
  \ifdefstring{#2}{是}{\vmark[#1]{\DeclYesX,#3}}{%
  \ifdefstring{#2}{否}{\vmark[#1]{\DeclNoX,#3}}{%
  \ifdefstring{#2}{是?}{\uvmark[#1]{\DeclYesX,#3}}{%
  \ifdefstring{#2}{否?}{\uvmark[#1]{\DeclNoX,#3}}{%
  \ifdefstring{#2}{是？}{\uvmark[#1]{\DeclYesX,#3}}{%
  \ifdefstring{#2}{否？}{\uvmark[#1]{\DeclNoX,#3}}{%
  \ifdraft\put(\DeclYesX,#3){\makebox[\dimexpr\DeclNoX bp-\DeclYesX bp\relax]{%
    \fontsize{9}{9}\selectfont\textcolor{red}{〔待定〕}}}\fi}}}}}}}

% \digit{\巨集}{第幾個字}{x,y}{字級}：把字串第 n 個字置中印在格子裡（中文大寫金額、郵遞區號）
\newcommand{\digit}[4]{\putc[#4]{#3}{%
  \ifunconf{#1}{\StrGobbleRight{#1}{1}[\digitstr]\unc{\StrChar{\digitstr}{#2}}}{\StrChar{#1}{#2}}}}

% ---------------------------------------------------------------- 開關
% \repname[字級]{x,y}{姓名}：\PrintRepOnSign = 1 才在簽章處印負責人姓名
\providecommand{\PrintRepOnSign}{0}
\newcommand{\repname}[3][12]{\ifdefstring{\PrintRepOnSign}{1}{\putl[#1]{#2}{#3}}{}}
% \ifsigndate{內容}：\PrintSignDate = 1 才印簽署日期（0 = 留白手寫）
\providecommand{\PrintSignDate}{1}
\newcommand{\ifsigndate}[1]{\ifdefstring{\PrintSignDate}{1}{#1}{}}
  % 再載入本檔
%  然後用 \includepdf[..., picturecommand*={\你的Overlay}]{src/表單.pdf}
%
%  值的三種狀態（填在 vendor.tex）：
%      {}        空白      → 草稿印紅色〔提示〕；正式版不印
%      {否?}     結尾有 ?  → 建議值、尚未確認：草稿印紅字；預覽版印黑字；正式版不印
%      {否}      其他      → 已確認：一律印黑字
%  模式：\drafttrue（草稿）／\draftfalse（正式）；preview.tex 會設 \PreviewMode（預覽版）。
% ============================================================================
\usepackage{fontspec}
\usepackage{xeCJK}
\usepackage{xcolor}
\usepackage{graphicx}
\usepackage{etoolbox}
\usepackage{xstring}
\expandarg
\usepackage{pdfpages}

% ---------------------------------------------------------------- 字型
% 依序找得到哪個就用哪個；要指定可在 vendor.tex 裡先 \def\TenderCJKFont{字型名稱}
\providecommand{\TenderCJKFont}{}
\providecommand{\TenderLatinFont}{}
\newcommand{\tryCJK}[1]{\ifdefempty{\TenderCJKFont}{\IfFontExistsTF{#1}{\def\TenderCJKFont{#1}}{}}{}}
\tryCJK{DFKai-SB}               % Windows 標楷體
\tryCJK{BiauKai}                % macOS 標楷體
\tryCJK{TW-Kai}                 % 全字庫正楷體
\tryCJK{Noto Serif CJK TC}
\tryCJK{Noto Sans CJK TC}
\ifdefempty{\TenderCJKFont}{\PackageError{overlay-core}{找不到中文字型，請在 vendor.tex 設定 \string\TenderCJKFont}{}}{}
\setCJKmainfont{\TenderCJKFont}
\ifdefempty{\TenderLatinFont}{%
  \IfFontExistsTF{Times New Roman}{\setmainfont{Times New Roman}}{\setmainfont{TeX Gyre Termes}}}%
  {\setmainfont{\TenderLatinFont}}

% \ifdraft、\ifpreview 由 forms.tex 在載入 vendor.tex 之前用 \newif 宣告
\ifdefined\PreviewMode\draftfalse\previewtrue\fi
\setlength{\unitlength}{1bp}

% ---------------------------------------------------------------- 放字
% \putl / \putc / \putr [字級]{x,y}{內容}：靠左／置中／靠右，y 為基線
\newcommand{\putl}[3][12]{\put(#2){\fontsize{#1}{#1}\selectfont #3}}
\newcommand{\putc}[3][12]{\put(#2){\makebox[0pt]{\fontsize{#1}{#1}\selectfont #3}}}
\newcommand{\putr}[3][12]{\put(#2){\makebox[0pt][r]{\fontsize{#1}{#1}\selectfont #3}}}

% ---------------------------------------------------------------- 確認狀態
% \ifunconf{值}{是}{否}：值是否以 ? 或 ？ 結尾
\newcommand{\ifunconf}[3]{\IfEndWith{#1}{?}{#2}{\IfEndWith{#1}{？}{#2}{#3}}}
% \unc{文字}：尚未確認的值——草稿紅字、預覽版黑字、正式版不印
\newcommand{\unc}[1]{\ifdraft\textcolor{red}{#1}\fi\ifpreview#1\fi}
% \val{\巨集}{提示}：必填欄位（空白時草稿印提示）
\newcommand{\val}[2]{\ifdefempty{#1}{\ifdraft\textcolor{red}{〔#2〕}\fi}%
  {\ifunconf{#1}{\unc{\StrGobbleRight{#1}{1}}}{#1}}}
% \opt{\巨集}：選填欄位（空白時不印）
\newcommand{\opt}[1]{\ifdefempty{#1}{}{\ifunconf{#1}{\unc{\StrGobbleRight{#1}{1}}}{#1}}}
% \draftnote{y}：頁首狀態說明（草稿紅字／預覽版灰字）
\newcommand{\draftnote}[1]{%
  \ifdraft\putl[8]{28,#1}{\textcolor{red}{草稿｜紅字＝尚未填寫或尚未確認。在 vendor.tex 補資料；值後面的「?」刪掉即視為確認。}}\fi
  \ifpreview\putl[8]{28,#1}{\textcolor{gray}{預覽版｜帶「?」的建議值也以黑字顯示，僅供檢查版面，勿用於投標。}}\fi}

% ---------------------------------------------------------------- 版面輔助
% \fitw{最大寬度}{內容}：超過寬度就等比縮小（長公司名、地址）
\newsavebox{\fitbox}
\newcommand{\fitw}[2]{\sbox{\fitbox}{#2}%
  \ifdim\wd\fitbox>#1\relax\resizebox{#1}{!}{\usebox{\fitbox}}\else\usebox{\fitbox}\fi}
% \putwrap[字級]{x,y}{寬度}{最大高度}{內容}：以 x 置中；放不下自動換行，再放不下等比縮小
\newsavebox{\wrapA}\newsavebox{\wrapB}
\newcommand{\putwrap}[5][14]{%
  \sbox{\wrapA}{\fontsize{#1}{#1}\selectfont #5}%
  \ifdim\wd\wrapA>#3\relax
    \sbox{\wrapA}{\parbox{#3}{\centering\fontsize{#1}{\dimexpr#1pt*23/20\relax}\selectfont #5}}%
    \sbox{\wrapB}{\ifdim\dimexpr\ht\wrapA+\dp\wrapA\relax>#4\relax\resizebox*{!}{#4}{\usebox{\wrapA}}\else\usebox{\wrapA}\fi}%
    \put(#2){\makebox[0pt]{\raisebox{\dimexpr#1pt*35/100-\ht\wrapB/2+\dp\wrapB/2\relax}{\usebox{\wrapB}}}}%
  \else
    \put(#2){\makebox[0pt]{\usebox{\wrapA}}}%
  \fi}

% ---------------------------------------------------------------- 勾選
\newcommand{\vmark}[2][13]{\put(#2){\makebox[0pt]{\fontsize{#1}{#1}\selectfont\bfseries V}}}
\newcommand{\rvmark}[2][13]{\put(#2){\makebox[0pt]{\fontsize{#1}{#1}\selectfont\bfseries\textcolor{red}{V}}}}
\newcommand{\uvmark}[2][13]{\ifdraft\rvmark[#1]{#2}\fi\ifpreview\vmark[#1]{#2}\fi}
% \chk{\開關}{x,y}：1 = 勾；1? = 尚未確認的勾；其他不勾
\newcommand{\chk}[2]{\ifdefstring{#1}{1}{\vmark[11]{#2}}{%
  \ifdefstring{#1}{1?}{\uvmark[11]{#2}}{\ifdefstring{#1}{1？}{\uvmark[11]{#2}}{}}}}

% \decl[V 字級]{\巨集}{y}：是／否兩欄的勾選（聲明書這類表格）。
%   欄位中心 x 由 \DeclYesX、\DeclNoX 設定（每份表單量一次）。
\providecommand{\DeclYesX}{474}
\providecommand{\DeclNoX}{527.5}
\newcommand{\decl}[3][13]{%
  \ifdefstring{#2}{是}{\vmark[#1]{\DeclYesX,#3}}{%
  \ifdefstring{#2}{否}{\vmark[#1]{\DeclNoX,#3}}{%
  \ifdefstring{#2}{是?}{\uvmark[#1]{\DeclYesX,#3}}{%
  \ifdefstring{#2}{否?}{\uvmark[#1]{\DeclNoX,#3}}{%
  \ifdefstring{#2}{是？}{\uvmark[#1]{\DeclYesX,#3}}{%
  \ifdefstring{#2}{否？}{\uvmark[#1]{\DeclNoX,#3}}{%
  \ifdraft\put(\DeclYesX,#3){\makebox[\dimexpr\DeclNoX bp-\DeclYesX bp\relax]{%
    \fontsize{9}{9}\selectfont\textcolor{red}{〔待定〕}}}\fi}}}}}}}

% \digit{\巨集}{第幾個字}{x,y}{字級}：把字串第 n 個字置中印在格子裡（中文大寫金額、郵遞區號）
\newcommand{\digit}[4]{\putc[#4]{#3}{%
  \ifunconf{#1}{\StrGobbleRight{#1}{1}[\digitstr]\unc{\StrChar{\digitstr}{#2}}}{\StrChar{#1}{#2}}}}

% ---------------------------------------------------------------- 開關
% \repname[字級]{x,y}{姓名}：\PrintRepOnSign = 1 才在簽章處印負責人姓名
\providecommand{\PrintRepOnSign}{0}
\newcommand{\repname}[3][12]{\ifdefstring{\PrintRepOnSign}{1}{\putl[#1]{#2}{#3}}{}}
% \ifsigndate{內容}：\PrintSignDate = 1 才印簽署日期（0 = 留白手寫）
\providecommand{\PrintSignDate}{1}
\newcommand{\ifsigndate}[1]{\ifdefstring{\PrintSignDate}{1}{#1}{}}
